% Options for packages loaded elsewhere
\PassOptionsToPackage{unicode}{hyperref}
\PassOptionsToPackage{hyphens}{url}
\documentclass[
  ukrainian,
]{article}
\usepackage{xcolor}
\usepackage[margin=2.5cm]{geometry}
\usepackage{amsmath,amssymb}
\setcounter{secnumdepth}{-\maxdimen} % remove section numbering
\usepackage{iftex}
\ifPDFTeX
  \usepackage[T1]{fontenc}
  \usepackage[utf8]{inputenc}
  \usepackage{textcomp} % provide euro and other symbols
\else % if luatex or xetex
  \usepackage{unicode-math} % this also loads fontspec
  \defaultfontfeatures{Scale=MatchLowercase}
  \defaultfontfeatures[\rmfamily]{Ligatures=TeX,Scale=1}
\fi
\usepackage{lmodern}
\ifPDFTeX\else
\setmainfont{DejaVu Serif}
\setsansfont{DejaVu Sans}
\setmonofont{DejaVu Sans Mono}
\setmathfont{STIX Two Math}
\fi
\ifPDFTeX\else
  % xetex/luatex font selection
\fi
% Use upquote if available, for straight quotes in verbatim environments
\IfFileExists{upquote.sty}{\usepackage{upquote}}{}
\IfFileExists{microtype.sty}{% use microtype if available
  \usepackage[]{microtype}
  \UseMicrotypeSet[protrusion]{basicmath} % disable protrusion for tt fonts
}{}
\makeatletter
\@ifundefined{KOMAClassName}{% if non-KOMA class
  \IfFileExists{parskip.sty}{%
    \usepackage{parskip}
  }{% else
    \setlength{\parindent}{0pt}
    \setlength{\parskip}{6pt plus 2pt minus 1pt}}
}{% if KOMA class
  \KOMAoptions{parskip=half}}
\makeatother
\usepackage{graphicx}
\graphicspath{{/Users/kalinicenkovitalijmikolajovic/.codex/worktrees/5c15/Script_for_translate_book/Output_Final/}}
\setkeys{Gin}{width=0.74\linewidth,height=0.78\textheight,keepaspectratio}
\makeatletter
\newsavebox\pandoc@box
\newcommand*\pandocbounded[1]{% scales image to fit in text height/width
  \sbox\pandoc@box{#1}%
  \Gscale@div\@tempa{\textheight}{\dimexpr\ht\pandoc@box+\dp\pandoc@box\relax}%
  \Gscale@div\@tempb{\linewidth}{\wd\pandoc@box}%
  \ifdim\@tempb\p@<\@tempa\p@\let\@tempa\@tempb\fi% select the smaller of both
  \ifdim\@tempa\p@<\p@\scalebox{\@tempa}{\usebox\pandoc@box}%
  \else\usebox{\pandoc@box}%
  \fi%
}
% Set default figure placement to htbp
\def\fps@figure{htbp}
\makeatother
\ifLuaTeX
\usepackage[bidi=basic,shorthands=off,provide=*]{babel}
\else
\usepackage[bidi=default,shorthands=off,provide=*]{babel}
\fi
\ifLuaTeX
  \usepackage{selnolig} % disable illegal ligatures
\fi
\setlength{\emergencystretch}{3em} % prevent overfull lines
\providecommand{\tightlist}{%
  \setlength{\itemsep}{0pt}\setlength{\parskip}{0pt}}
\usepackage{bookmark}
\IfFileExists{xurl.sty}{\usepackage{xurl}}{} % add URL line breaks if available
\urlstyle{same}
\hypersetup{
  pdflang={uk},
  hidelinks,
  pdfcreator={LaTeX via pandoc}}

\author{}
\date{}

\begin{document}

\pandocbounded{\includegraphics[keepaspectratio]{images/_page_315_Figure_7.jpeg}}

Фіг. 42 2. Обчислення інтеграції Монте-Карло.

Очевидно, якщо випадковим чином кинути близько 100 дротиків у
прямокутник R, що містить графік функції, можна отримати оцінку значення
інтеграла, помноживши площу прямокутника R на відносну кількість
дротиків під графіком функції. Отже, результат нашої гри

т. е. получается УЧП с переменными коэффициентами.

\begin{verbatim}
3. Оскільки хвильове рівняння  u_{tt} = \alpha^2 u_{xx}  містить похідну другого порядку часу, щоб отримати єдине розв'язання в точці t>0, необхідно вказати  \partial sa  початкові умови:
\end{verbatim}

--- це інтегральний бал 1.

Щоб виконати ці обчислення на комп'ютері, нам потрібно якимось чином
отримати послідовність випадкових точок (подивимось, як це буде зроблено
пізніше), тобто наказати комп'ютеру «кидати дротики». Блок-схема на рис.
42.3 показує, як комп'ютер вирішить цю проблему.

\pandocbounded{\includegraphics[keepaspectratio]{images/_page_316_Figure_2.jpeg}}

\textbf{РИС.} 42.3. Блок-схема розрахунку інтегрального
\texttt{\textbackslash{}int\_\{a\}\^{}\{b\}\ f(x)\ dx} методом
Монте-Карло (100 кидків);
\texttt{M\ =\ \textbackslash{}max\_\{a\}\ f(x)} .

\section{Випадкові
числа}\label{ux432ux438ux43fux430ux434ux43aux43eux432ux456-ux447ux438ux441ux43bux430}

Перш ніж перейти до рівняння Монте-Карло з частковими похідними,
зупинимося на важливому питанні випадкових чисел. Щойно, обчислюючи
інтеграл, нам потрібно було побудувати послідовність випадкових точок
\texttt{P\_i\ =\ (x\_i,\ y\_i)} всередині прямокутника R. Іншими
словами, x-координата точки P має бути випадковим числом з інтервалу
{[}a, b{]}, а \texttt{y\_i} від інтервалу {[}0, M{]}. Щоб знайти
випадкові числа з цих інтервалів, звертаємося до послідовності
випадкових чисел, рівномірно розподілених по сегменту {[}0, 1{]}. Тоді
числа, рівномірно розподілені на сегменті {[}a, b{]}, можна обчислити за
формулою

(початкове зміщення струни), \texttt{u\_t}(x, 0) = \texttt{g(x)}
(початкова швидкість струни).

Таким чином, хвильове рівняння відрізняється від рівняння
теплопровідності, для якого потрібно встановити лише одну початкову
умову.

\begin{enumerate}
\def\labelenumi{\arabic{enumi}.}
\setcounter{enumi}{3}
\tightlist
\item
  За допомогою хвильового рівняння можна описати розподіл електричного
  струму в дроті. З законів Кірхгофа отримуємо таку систему двох рівнянь
  у похідних похідних першого порядку:
\end{enumerate}

(16.3)

Тепер, звичайно, виникає питання: як отримати послідовність випадкових
чисел,
\texttt{\textbackslash{}\{r\_i:\ i=1,2,\textbackslash{}ldots\textbackslash{}\}}
рівномірно розподілених на {[}0,1{]}? Найпоширенішим методом сьогодні є
метод дедукцій (метод порівнянь). Використовуючи цей метод, отримуємо
послідовність випадкових цілих чисел (наприклад, 2120, 1401, 177,
3013,\ldots), потім ліворуч від кожного числа розміщується десяткова
крапка, так що всі числа знаходяться між нулем і одним
\texttt{(0,212;\ 0,1401;\ 0,0177;\ 0,3013;\ \textbackslash{}ldots)} .

Отже, щоб отримати послідовність випадкових цілих чисел (між 0 і P), ми
використаємо метод віднімання.

\section{Отримання випадкових чисел за допомогою методу
відрахування.}\label{ux43eux442ux440ux438ux43cux430ux43dux43dux44f-ux432ux438ux43fux430ux434ux43aux43eux432ux438ux445-ux447ux438ux441ux435ux43b-ux437ux430-ux434ux43eux43fux43eux43cux43eux433ux43eux44e-ux43cux435ux442ux43eux434ux443-ux432ux456ux434ux440ux430ux445ux443ux432ux430ux43dux43dux44f.}

\begin{enumerate}
\def\labelenumi{\arabic{enumi}.}
\item
  Як перше випадкове ціле число, обирайте будь-яке число, що знаходиться
  між 0 і P (число P обирається наперед).
\item
  Помножте це випадкове число на певне фіксоване число.
\end{enumerate}

Попередньо обраний множник M.

\begin{enumerate}
\def\labelenumi{\arabic{enumi}.}
\setcounter{enumi}{2}
\item
  Додайте до добутку деяке фіксоване ціле число \emph{K}, обране
  заздалегідь.
\item
  Поділити отриману суму на P і залишок
\end{enumerate}

Виберіть нове випадкове число.

Метод виведення можна записати у формі формули
\texttt{r\_\{i+1\}\ \textbackslash{}equiv\ (Mr\_i\ +\ K)\ \textbackslash{}bmod\ P}
, що означає, що для отримання наступного \texttt{r\_\{i+1\}} потрібно
взяти попередній \texttt{r\_i} , помножити його на M, додати результат
на K, поділити суму на P і взяти решту поділу. Наприклад, якщо P = 100,
M = 37, K = 16, \texttt{r\_0\ =\ 15} , то отримується наступна
послідовність випадкових чисел:

где

ж --- координата вдоль провода,

\emph{t} --- время,

\begin{verbatim}
i(x, t) — распределение тока вдоль провода,  v\left(x,\ t\right)  — распределение потенциала вдоль провода,
\end{verbatim}

С -- емкость провода на единицу длины,

G --- утечка на единицу длины,

Р --- сопротивление на единицу длины,

-- индуктивность провода на единицу длины.

Рівняння (16.3) зазвичай називають системою телеграфних рівнянь. У такій
формі вони зазвичай не використовуються. Щоб перетворити цю систему,
диференціюйте перше рівняння на x, друге на t, помножте обидві частини
на C і відніміть від першого. У результаті ми отримуємо

\begin{verbatim}
Mr_1 + K = 2643
\end{verbatim}

Воспользуемся теперь вторым уравнением из (16.3)

\begin{verbatim}
r = 71
\end{verbatim}

окончательно получаем

(16.4)

Y

Це рівняння для струму, відоме як телеграфне рівняння, є гіперболічним
рівнянням другого порядку (якщо, звісно, C і L не дорівнюють нулю,
інакше рівняння стає параболічним).

Напряжение описывается точно таким же уравнением:

(16.5)

\begin{verbatim}
I = \int_{0}^{1} \int_{0}^{1} \int_{0}^{1} \int_{0}^{1} e^{-(x^{2}+y^{2}+z^{2}+\omega^{2})} dx dy dz dw,
\end{verbatim}

Если G = R = 0, то уравнения (16.4)---(16.5) упрощаются:

(16.6)

\begin{verbatim}
I = \int_{0}^{1} e^{\sin x} dx.
\end{verbatim}

\section{\texorpdfstring{ЗАВДАННЯ - 1. Отримайте рівняння (16.5) для
\emph{and} з системи з двох рівнянь першого порядку
(16.3)}{ЗАВДАННЯ - 1. Отримайте рівняння (16.5) для and з системи з двох рівнянь першого порядку (16.3)}}\label{ux437ux430ux432ux434ux430ux43dux43dux44f---1.-ux43eux442ux440ux438ux43cux430ux439ux442ux435-ux440ux456ux432ux43dux44fux43dux43dux44f-16.5-ux434ux43bux44f-and-ux437-ux441ux438ux441ux442ux435ux43cux438-ux437-ux434ux432ux43eux445-ux440ux456ux432ux43dux44fux43dux44c-ux43fux435ux440ux448ux43eux433ux43e-ux43fux43eux440ux44fux434ux43aux443-16.3}

\begin{itemize}
\tightlist
\item
  \begin{enumerate}
  \def\labelenumi{\arabic{enumi}.}
  \setcounter{enumi}{1}
  \tightlist
  \item
    Знаючи фізичне значення кожного члена хвильового рівняння, що ви
    можете сказати про поведінку в момент розв'язання наступної задачі:
  \end{enumerate}
\end{itemize}

(YYII)

\begin{verbatim}
r_{i+1} \equiv (3r_i + 4) \mod 7, \quad r_0 = 0.
,  0 < x < 1  ,  0 < t < \infty  ,
\end{verbatim}

(ГУ)

\begin{verbatim}
I = \int_{0}^{1} \int_{0}^{1} \int_{0}^{1} e^{-(x^{2}+y^{2}+z^{2})} dx dy dz.
u(x,y) = g(x,y) =
\begin{cases} u(0, t) = 0, & 0 < t < \infty, \\ u(1, t) = 0, & 0 < t < \infty, \end{cases}
\end{verbatim}

Після того, як цю конкретну задачу буде розв'язано методом Монте-Карло,
ми розглянемо, як розв'язувати більш загальні задачі за допомогою цього
методу.

Щоб проілюструвати метод Монте-Карло, розглянемо гру під назвою
«Блукаючий п'яниця». Щоб грати в цю гру, потрібна таблиця з сіткою
ліній, як показано на рис. 43.1.

Перейдем к описанию правил игры.

\section{Как играть в «Блуждающего
пьяницу»?}\label{ux43aux430ux43a-ux438ux433ux440ux430ux442ux44c-ux432-ux431ux43bux443ux436ux434ux430ux44eux449ux435ux433ux43e-ux43fux44cux44fux43dux438ux446ux443}

КРОК 1. Блукання п'яниці починаються з довільної точки сітки (у нашому
випадку --- точки А).

\pandocbounded{\includegraphics[keepaspectratio]{images/_page_321_Figure_4.jpeg}}

\begin{itemize}
\item
  Початкова точка - - внутрішні точки сітки - □ - Кінцева точка p;
\item
  \texttt{g\_i} винагороди в m. \texttt{p\_I} Фіг. 43.1. Ігрове поле для
  гри «Блукаючий п'яниця».
\end{itemize}

КРОК 2. У кожному ході гри п'яний випадково переходить до однієї з
чотирьох суміжних точок сітки (у нашому випадку сусідні точки для A ---
це точки B, C, D і E, і ймовірність досягти кожної з цих точок становить
0,25).

Таблица 43.1

\begin{verbatim}
Вероятность окончания случайного блуждания в граничной точке p_i и величина вознаграждения g_i | Граничная точка р <sub>і</sub>                        | P_{\mathcal{A}}\left(p_{i}\right) — относительная доля блужданий, закончившихся в точке p_{i} | g_i — величина вознаграждения при достижении точки p_i |  |  |
\end{verbatim}

\textbar-------------------------------------------------------\textbar---------------------------------------------------------------------------------------------------\textbar------------------------------------------------------------\textbar--\textbar--\textbar{}

1234567891011 \textbar{}
0,040,150,030,060,170,050,060,150,030,060,160,04 \textbar{} 110000000
\textbar{} \textbar{} \textbar{}

ti

КРОК 3. Після переходу до наступної точки процес блукання відновлюється.
Таким чином, п'яниця ходить від точки до точки, поки випадково не
опиняється на межі \texttt{p\_t} . Тут це зупиняється, і ми записуємо
номер точки \texttt{p\_t} . Це кінець однієї невимушеної прогулянки.

КРОК 4. Повторимо 1-3 рази достатньо і знайдемо відносну кількість
прибуттів п'яниця в кожній з меж \texttt{p\_i} . Таблиця 1. 43.1 показує
типові результати, отримані внаслідок 100 випадкових прогулянок.

\begin{verbatim}
КРОК B: Припустимо, що п'яниця отримує винагороду g_i = g\left(p_i\right) (значення граничної умови в точці p_i , якщо він досягає точки p_i ) і припустимо, що мета гри — обчислити середній R\left(A\right) винагороди за всі випадкові прогулянки, починаючи з точки A. Цей середній виграш визначається формулою
\end{verbatim}

\texttt{Гра\ закінчується,\ коли\ знайдено\ значення\ \textquotesingle{}R\textbackslash{}left(A\textbackslash{}right)\textquotesingle{}\ .\ Відповідно\ до\ даних\ таблиці.\ 43.1\ за\ \textquotesingle{}R\textbackslash{}left(A\textbackslash{}right)\textquotesingle{}\ отримуємо}

Теперь объясним, зачем нам понадобилось играть в такую игру.

\section{Для чего играть в «Блуждающего
пьяницу»?}\label{ux434ux43bux44f-ux447ux435ux433ux43e-ux438ux433ux440ux430ux442ux44c-ux432-ux431ux43bux443ux436ux434ux430ux44eux449ux435ux433ux43e-ux43fux44cux44fux43dux438ux446ux443}

Виявляється, що середня винагорода, про яку щойно йдеться, є приблизним
розв'язком проблеми Діріхле в точці A. Це цікаве спостереження базується
на двох фактах.

\begin{enumerate}
\def\labelenumi{\arabic{enumi}.}
\item
  Припустимо, що п'яниця почав свій шлях з точки, що лежить на Єрениці.
  Кожна така прогулянка одразу закінчується в одному й тому ж місці, і
  птах отримує винагороду \texttt{g\_i} . Отже, середня винагорода за
  кожну граничну точку дорівнює \texttt{g\_i} .
\item
  Тепер припустимо, що прогулянка починається з внутрішньої точки. Тоді
  стає зрозуміло, що середня винагорода за бал \texttt{R(A)} буде
  арифметичним середнім середніх нагород за чотири сусідні бали
\end{enumerate}

\begin{verbatim}
\begin{cases} u(x, 0) = \sin(\pi x), & 0 \le x \le 1, \\ u_t(x, 0) = 0. & 0 \le x \le 1, \end{cases}
\end{verbatim}

Давайте знову запитаємо себе: чому середнє значення винагороди
\texttt{R(A)} наближається до розв'язання задачі Діріхле в точці A? Ми
бачили, що значення \texttt{R(A)} задовольняє двох

уравнениям

(во внутренних точках), \texttt{R(A)\ =\ g\_i} (в граничных точках).

Якщо \texttt{g\_i} --- значення функції \texttt{g(x,\ y)} з граничної
умови у граничних точках \texttt{p\_i} , то наші два рівняння точно
збігаються з двома рівняннями, отриманими при розв'язанні задачі Діріхле
методом скінченної різниці. Тобто значення \texttt{R(A)} відповідає
значенню \texttt{u\_\{i,f\}} у різницевих рівняннях

(i, j)---внутренняя точка,
\texttt{u\_\{ij\}\ =\ g\_\{ij\},\ g\_\{ij\}}---значение решения в
граничной точке (i, j).

Отже, значення \texttt{R(A)} дійсно наближає розв'язок рівняння
часткових у точці A.

Тепер можна вважати, що гра «Блукаючий п'яниця» складається з трьох
кроків.

\section{Розв'язання рівняння Лапласа за методом
Монте-Карло}\label{ux440ux43eux437ux432ux44fux437ux430ux43dux43dux44f-ux440ux456ux432ux43dux44fux43dux43dux44f-ux43bux430ux43fux43bux430ux441ux430-ux437ux430-ux43cux435ux442ux43eux434ux43eux43c-ux43cux43eux43dux442ux435-ux43aux430ux440ux43bux43e}

\begin{itemize}
\item
  Наступні правила дозволяють отримати розв'язок в одній внутрішній
  точці квадрата.
\item
  КРОК 1. Давайте зробимо кілька випадкових прогулянок, починаючи з
  точки A і закінчуючи однією з прикордонних точок. Давайте
  відстежувати, скільки разів прогулянки закінчувалися на кожній межі.
\item
  КРОК 2. Після завершення всіх прогулянок для кожної граничної точки ми
  обчислюємо відносну кількість прогулянок, які закінчилися в цій точці.
  Позначимо ці значення \texttt{P\_A}(\texttt{p\_i}) .
\end{itemize}

\begin{verbatim}
- КРОК 3. Обчислюємо наближений розв'язок u(A) за формулою u(A) = g_1 P_A(p_1) + g_2 P_A(p_2) + \dots + g_N P_A(p_N) ,
\end{verbatim}

де \texttt{g\_i} --- значення функції g у точці \texttt{p\_i} , а N ---
кількість граничних точок.

Блукаючий п'яниця можна покращити, щоб розв'язати складнішу задачу.
Нижче наведено приклад такого оновлення.

Розв'язок задачі Діріхле з змінними коефіцієнтами.

Рассмотрим эллиптическую краевую вадачу в квадрате

Дайте физическую интерпретацию этой задачи.

\begin{enumerate}
\def\labelenumi{\arabic{enumi}.}
\setcounter{enumi}{3}
\tightlist
\item
  Для уравнения \texttt{u\_\{tt\}\ =\ u\_\{xx\}} найдите все решения
  вида
\end{enumerate}

\begin{verbatim}
u_{xx} + (\sin x) u_{yy} = 0 , 0 < x < \pi , 0 < y < \pi , (\Gamma Y) u(x, y) = g(x, y) на границе квадрата.
\end{verbatim}

\begin{verbatim}
Щоб розв'язати цю задачу, замінимо u_{xx} , u_{yy} та \sin x наступним чином:
\end{verbatim}

Будет ли решением сумма двух решений такого вида?

\section{МЕТАЛЕКЦІЇ ФОРМУЛИ Д'АЛАМБЕРА: Побудуйте розв'язок проблеми
Коші для хвильового
рівняння}\label{ux43cux435ux442ux430ux43bux435ux43aux446ux456ux457-ux444ux43eux440ux43cux443ux43bux438-ux434ux430ux43bux430ux43cux431ux435ux440ux430-ux43fux43eux431ux443ux434ux443ux439ux442ux435-ux440ux43eux437ux432ux44fux437ux43eux43a-ux43fux440ux43eux431ux43bux435ux43cux438-ux43aux43eux448ux456-ux434ux43bux44f-ux445ux432ux438ux43bux44cux43eux432ux43eux433ux43e-ux440ux456ux432ux43dux44fux43dux43dux44f}

и
\emph{\{yy\}=\texttt{u\_\{l+1,\textbackslash{}\ j\}}-2u}\{l,~j\}+\texttt{u\_\{l-1,\textbackslash{}\ j\}},
(центральные разностные производные),

і підставляють їх у рівняння похідних похідних. Відповідна сітка
показана на рис. 43.2. Розв'язавши отримане рівняння за
\texttt{u\_\{i,l\}} , отримуємо

Давайте розглянемо це рівняння уважніше. Коефіцієнти при
\texttt{u\_\{i+1,\textbackslash{}\ j\},\textbackslash{}\ u\_i,\textbackslash{}\ j+i,\textbackslash{}\ u\_l,\textbackslash{}\ j-1}
n \texttt{u\_\{i-1\},\textbackslash{}\ j} додатні, і їхня сума дорівнює
одиниці. Іншими словами, \texttt{u\_i,\textbackslash{}\ j} рішення є
зваженим середнім рішень у чотирьох сусідніх точках. Отже, можливо P і C

\pandocbounded{\includegraphics[keepaspectratio]{images/_page_324_Picture_9.jpeg}}

. 43.2. Узлы сетки для случайного маршруга.

модифікувати Блукаючого П'яницю так, щоб ймовірності переміщення до
сусідніх точок не були 0,25, а дорівнювали коефіцієнту відповідного
члена. Інакше кажучи, якщо п'яниця знаходиться в точці (i, j), то він
переходить у точку

Ця задача описує рух нескінченної струни за заданих початкових умов.
Його розв'язав у 1750 році французький математик д'Аламбер. Формула
Д'Аламбера

с вероятностью
\texttt{\textbackslash{}frac\{1\}\{2(1+\textbackslash{}sin\ x\_j)\}} ,
\texttt{(i,\ j-1)} с вероятностью
\texttt{\textbackslash{}frac\{1\}\{2(1+\textbackslash{}sin\ x\_j)\}} ,
\texttt{(i+1,\ j)} с вероятностью
\texttt{\textbackslash{}frac\{\textbackslash{}sin\ x\_j\}\{2(1+\textbackslash{}sin\ x\_j)\}}
, \texttt{(i-1,\ j)} с вероятностью
\texttt{\textbackslash{}frac\{\textbackslash{}sin\ x\_j\}\{2(1+\textbackslash{}sin\ x\_j)\}}
.

Інакше гра не змінюється. Модифікація для інших завдань може бути більш
винахідливою, але основна ідея залишається незмінною. Читач може почати
розробляти власні ігри, щоб розв'язувати інші проблеми. Параболічний
випадок розглядається в розділі «Проблеми». \# ПОВАГА
\texttt{1.\ Зверніть\ увагу,\ що\ якщо\ відносні\ частоти\ удару\ \textquotesingle{}P\_A(p\_l)\textquotesingle{}\ відомі,\ то\ значення\ \textquotesingle{}u\textbackslash{}left(A\textbackslash{}right)\textquotesingle{}\ за\ інших\ граничних\ умов\ \textquotesingle{}g\_l\textquotesingle{}\ легко\ знайти.\ Для\ цього\ достатньо\ підставити\ \textquotesingle{}g\_l\textquotesingle{}\ у\ формулу}

полегшує пошук розв'язку, якщо початкові умови відомі. Крім того, це
дозволяє дати цікаву фізичну інтерпретацію розв'язку мовою хвиль, що
рухаються.

\begin{verbatim}
Як читач, ймовірно, пам'ятає, у параболічному випадку ми спочатку розв'язали задачі дифузії на скінченному відрізку (методом розділення змінних), а лише потім переходили до необмежених інтервалів (-  \infty < x < \infty  ), де застосовували перетворення Фур'є. У гіперболічному випадку ми зробимо навпаки. Почнемо з розв'язання одномірного хвильового рівняння на всій дійсній осі, тобто з задачі Коші
\end{verbatim}

(17.1)

а не разыгрывать вновь случайные прогулки.

\begin{enumerate}
\def\labelenumi{\arabic{enumi}.}
\setcounter{enumi}{1}
\tightlist
\item
  У багатьох випадках дослідник зацікавлений у розв'язанні
  диференціального рівняння в певній точці. Якщо межа області є
  комплексною і рівняння залежить від трьох або чотирьох пояснювальних
  змінних, метод Монте-Карло може бути зручним. Методи Монте-Карло
  виникли для розв'язання дуже складних задач дифузії нейтронів. Ці
  завдання не можна розв'язати аналітично. \# ЗАВДАННЯ
\item
  Побудуйте блок-схему гри «Блукаючий п'яниця», щоб розв'язати проблему
\end{enumerate}

(УЧП)

, 0 \textless{} x \textless{} 1 , 0 \textless{} y \textless{} 1 , (ГУ)
\texttt{u(x,\ y)} = \texttt{g(x,\ y)} на границе,

\begin{itemize}
\item
  у вузлах внутрішньої сітки, якщо кількість вузлів сітки довільна. 2.
  Напишіть програму для розв'язання задачі відповідно до блок-схеми з
  завдання 1.
\item
  \begin{enumerate}
  \def\labelenumi{\arabic{enumi}.}
  \setcounter{enumi}{2}
  \tightlist
  \item
    Як модифікувати гру для вирішення проблеми
  \end{enumerate}
\end{itemize}

(УЧП)

містить лише початкові умови. Цю проблему можна розв'язати або
перетворенням Фур'є (для змінної t), або перетворенням Лапласа (для
змінної t). Однак ми використаємо інший метод (метод канонічних
координат), який, сподіваємося, зацікавить читача. Цей метод базується
на тих самих ідеях, що й метод переходу до рухомої системи відліку, з
яким ми познайомилися на лекції 15. Тож почнемо розв'язувати задачу
(17.1).

\section{Розв'язок одновимірного хвильового рівняння. Формула
Д'Аламбера}\label{ux440ux43eux437ux432ux44fux437ux43eux43a-ux43eux434ux43dux43eux432ux438ux43cux456ux440ux43dux43eux433ux43e-ux445ux432ux438ux43bux44cux43eux432ux43eux433ux43e-ux440ux456ux432ux43dux44fux43dux43dux44f.-ux444ux43eux440ux43cux443ux43bux430-ux434ux430ux43bux430ux43cux431ux435ux440ux430}

Розв'язок задачі (17.1) буде поділений на кілька етапів:

КРОК 1 (Замінити координати (x, t) на нові канонічні координати
\texttt{(\textbackslash{}xi,\ \textbackslash{}eta)} ).

Щоб розв'язати задачу (17.1), ми використаємо той факт, що якщо замінити
дві незалежні змінні x і t на нові координати простору-часу

, 0 \textless{} x \textless{} 1 , 0 \textless{} y \textless{} 1 , (ГУ)
\texttt{u(x,\ y)} = \texttt{g(x,\ y)} на границе?

Как выглядит в этой игре случайная прогулка?

\begin{enumerate}
\def\labelenumi{\arabic{enumi}.}
\setcounter{enumi}{3}
\tightlist
\item
  Можете ли вы так модифицировать игру, чтобы решить задачу
\end{enumerate}

\begin{verbatim}
\begin{array}{l} u_{xx} + u_{yy} = 0, \quad 0 < x < 1, \quad 0 < y < 1, \\ u(x, 1) = 0, \quad 0 < x < 1, \quad 0 < y < 1, \\ u(x, 0) = 0, \quad 0 < x < 1, \quad 0 < y < 1, \\ u(0, y) = 1, \quad \frac{\partial u}{\partial x}(1, y) = 0 \end{array}
\end{verbatim}

то уравнение

\begin{verbatim}
\begin{aligned}
(\text{УЧП})\quad & u_t = \alpha^2 u_{xx} - \beta u_x - \gamma u, && 0 < x < 1,\ 0 < t < \infty, \\
(\text{ГУ})\quad & u(0,t) = f(t),\quad u(1,t) = g(t), && 0 < t < \infty, \\
(\text{НУ})\quad & u(x,0) = \varphi(x), && 0 \leqslant x \leqslant 1.
\end{aligned}
\end{verbatim}

преобразуется в уравнение

ПРИМІТКА. Замінимо це рівняння на скінченну різницю Кренка--Ніколсона,
виразимо \texttt{u\_\{i+1,j\}} через п'ять значень
\texttt{u\_\{i+1,j-1\}} , \texttt{u\_\{i+1,j+1\}} ,
\texttt{u\_\{i,j+1\}} , \texttt{u\_\{i,j+1\}} , \texttt{u\_\{i,j+1\}} у
суміжних точках.

\section{КАЛЬКУЛЮС ВАРІАЦІЙ (РІВНЯННЯ
ЕЙЛЕРА-ЛАГРАНЖА)}\label{ux43aux430ux43bux44cux43aux443ux43bux44eux441-ux432ux430ux440ux456ux430ux446ux456ux439-ux440ux456ux432ux43dux44fux43dux43dux44f-ux435ux439ux43bux435ux440ux430-ux43bux430ux433ux440ux430ux43dux436ux430}

МЕТА ЛЕКЦІЇ: Ввести поняття функціоналу, тобто функції від функції, і
пояснити, як функціонали природно виникають у фізиці. Типовим
функціоналом є інтеграл

\begin{verbatim}
J[y] = \int_a^b F(x, y, y')\,dx,
\end{verbatim}

який є функцією від функції \texttt{y} (при цьому вважається, що
підінтегральна функція \texttt{F(x,\ y,\ y\textquotesingle{})} задана).
Приклад функціоналу:

\begin{verbatim}
J[y] = \int_0^1 \left[y^2(x) + y'^2(x)\right] dx.
\end{verbatim}

Ми покажемо, як знайти функцію
\texttt{\textbackslash{}overline\{y\}(x)}, що мінімізує функціонал
\texttt{J{[}y{]}}. Виявляється, що мінімізуюча функція
\texttt{\textbackslash{}overline\{y\}} повинна задовольняти так зване
рівняння Ейлера-Лагранжа. Воно відіграє ту саму роль, що й необхідна
умова мінімуму

\begin{verbatim}
\frac{df(x)}{dx} = 0
\end{verbatim}

для функції \texttt{f(x)} у точці \texttt{x} в диференціальному
численні.

Варіаційне числення тісно пов'язане з диференціальними рівняннями, але,
на жаль, більшість студентів його не вивчає. У цій лекції (і наступній)
подається вступ до варіаційного числення та показано, як можна
розв'язувати диференціальні рівняння в частинних похідних на основі
варіаційних принципів.

Варіаційне числення виникло одночасно з математичним аналізом у зв'язку
з розв'язанням задач максимізації та мінімізації функцій від функцій,
тобто функціоналів. Першою задачею варіаційного числення була задача про
брахістохрону, сформульована Йоганном Бернуллі у 1696 році. У цій задачі
потрібно знайти криву \texttt{y(x)} так, щоб мінімізувати час спуску без
тертя з однієї точки в іншу (див. рис. 44.1).

Бернуллі показав, що час спуску записується у вигляді

\begin{verbatim}
T = \int_0^T dt = \int_a^L \frac{dt}{ds}\,ds = \int_0^L \frac{ds}{v}
= \frac{1}{\sqrt{2mg}} \int_0^L \frac{ds}{\sqrt{y}}
= \frac{1}{\sqrt{2mg}} \int_a^b \sqrt{\frac{1+y'^2}{y}}\,dy
\end{verbatim}

і, отже, є функціоналом \texttt{T\textbackslash{}\{y\textbackslash{}\}}
від функції. Оскільки багато функціоналів мають подібну будову, ми
зосередимося на вивченні функціоналу загального вигляду

(44.1)

\begin{verbatim}
J[y] = \int_a^b F(x, y, y')\,dx.
\end{verbatim}

Тепер сформулюємо головну мету лекції: знайти функцію, яка доставляє
мінімум (або максимум) функціоналу (44.1). Стратегія пошуку буде такою
ж, як і при пошуку мінімуму функції у звичайному диференціальному
численні. Там ми шукали критичні точки з умови
\texttt{f\textquotesingle{}(x)=0}. У варіаційному численні все дещо
складніше, бо аргументом є не число, а функція. Проте загальний підхід
той самий: ми обчислимо функціональну похідну за функцією \texttt{y(x)}
і прирівняємо її до нуля. Нове рівняння буде аналогом умови
\texttt{df(x)/dx\ =\ 0}, але тепер це буде звичайне диференціальне
рівняння, відоме як рівняння Ейлера-Лагранжа.

Мінімізація функціоналу

\begin{verbatim}
J[y] = \int_a^b F(x, y, y')\,dx
\end{verbatim}

Розглянемо задачу знаходження функції \texttt{y(x)}, яка мінімізує цей
функціонал у класі гладких функцій, що задовольняють граничним умовам

\begin{verbatim}
y(a) = A, \qquad y(b) = B.
\end{verbatim}

(див. рис. 44.2).

\texttt{Нехай\ шукана\ функція\ \textquotesingle{}y\ =\ \textbackslash{}overline\{y\}(x)\textquotesingle{}\ існує,\ і\ розглянемо\ малу\ варіацію\ цієї\ функції,\ тобто\ функцію\ \textquotesingle{}\textbackslash{}overline\{y\}\ +\ \textbackslash{}varepsilon\textbackslash{}eta(x)\textquotesingle{},\ де\ \textquotesingle{}\textbackslash{}varepsilon\textquotesingle{}\ —\ мале\ число,\ а\ \textquotesingle{}\textbackslash{}eta(x)\textquotesingle{}\ —\ гладка\ функція,\ яка\ задовольняє\ граничним\ умовам\ \textquotesingle{}\textbackslash{}eta(a)=\textbackslash{}eta(b)=0\textquotesingle{}.}

\texttt{Тоді,\ якщо\ обчислити\ інтеграл\ \textquotesingle{}J\textquotesingle{}\ для\ близької\ функції\ \textquotesingle{}\textbackslash{}overline\{y\}\ +\ \textbackslash{}varepsilon\textbackslash{}eta\textquotesingle{},\ функціонал\ зросте,\ тобто}

\begin{verbatim}
J[\overline{y}] \leqslant J[\overline{y} + \varepsilon\eta]
\end{verbatim}

\texttt{для\ всіх\ \textquotesingle{}\textbackslash{}varepsilon\textquotesingle{}.\ Іншими\ словами,\ графік\ \textquotesingle{}\textbackslash{}varphi(\textbackslash{}varepsilon)=J{[}\textbackslash{}overline\{y\}+\textbackslash{}varepsilon\textbackslash{}eta{]}\textquotesingle{}\ як\ функції\ \textquotesingle{}\textbackslash{}varepsilon\textquotesingle{}\ має\ вигляд,\ показаний\ на\ рис.\ 44.3.}

\pandocbounded{\includegraphics[keepaspectratio]{images/_page_329_Figure_5.jpeg}}

Рис. 44.1. Задача про брахістохрону (з цієї задачі почалося варіаційне
числення).

\pandocbounded{\includegraphics[keepaspectratio]{images/_page_329_Figure_7.jpeg}}

\texttt{Рис.\ 44.2.\ Варіація\ функції:\ a\ —\ варіація\ функції\ \textquotesingle{}\textbackslash{}overline\{y\}(x)+\textbackslash{}varepsilon\textbackslash{}eta(x)\textquotesingle{};\ б\ —\ мінімізуюча\ функція\ \textquotesingle{}\textbackslash{}overline\{y\}(x)\textquotesingle{}.}

З рис. 44.3 видно, що тут потрібно обчислити похідну від

\begin{verbatim}
\varphi(\varepsilon) = J[\overline{y} + \varepsilon\eta]
\end{verbatim}

за \texttt{\textbackslash{}varepsilon}, покласти
\texttt{\textbackslash{}varepsilon\ =\ 0} і прирівняти отриманий вираз
до нуля, тобто

\begin{verbatim}
\frac{d\varphi(\varepsilon)}{d\varepsilon}
= \frac{d}{d\varepsilon} J[\overline{y} + \varepsilon\eta]\bigg|_{\varepsilon=0}
= \int_a^b \left[
\frac{\partial F}{\partial \overline{y}}\eta(x)
+ \frac{\partial F}{\partial \overline{y'}}\eta'(x)
\right] dx = 0.
\end{verbatim}

(Читач має виконати це обчислення самостійно.) Інтегруючи частинами,
дістаємо

\begin{verbatim}
\frac{d\varphi(\varepsilon)}{d\varepsilon}
\equiv \int_a^b \left\{
\frac{\partial F}{\partial \overline{y}} - \frac{d}{dx}\left[\frac{\partial F}{\partial \overline{y'}}\right]
\right\} \eta(x)\,dx = 0.
\end{verbatim}

Оскільки цей інтеграл дорівнює нулю для будь-якої функції
\texttt{\textbackslash{}eta(x)}, яка задовольняє граничним умовам
\texttt{\textbackslash{}eta(a)=\textbackslash{}eta(b)=0}, то
підінтегральний вираз має дорівнювати нулю, тобто

\pandocbounded{\includegraphics[keepaspectratio]{images/_page_330_Figure_6.jpeg}}

\texttt{Рис.\ 44.3.\ Графік\ функції\ \textquotesingle{}J{[}y+\textbackslash{}varepsilon\textbackslash{}eta{]}\textquotesingle{}\ в\ околі\ \textquotesingle{}\textbackslash{}varepsilon=0\textquotesingle{}.}

(44.2)

\begin{verbatim}
\frac{\partial F}{\partial \overline{y}} - \frac{d}{dx}\left[\frac{\partial F}{\partial \overline{y'}}\right] = 0
\end{verbatim}

(рівняння Ейлера-Лагранжа).

Рівняння (44.2) називається рівнянням Ейлера-Лагранжа. Хоча в загальному
вигляді воно здається складним, після підстановки конкретної функції
\texttt{F(x,y,y\textquotesingle{})} воно перетворюється на звичайне
диференціальне рівняння другого порядку відносно невідомої функції
\texttt{\textbackslash{}overline\{y\}(x)}. Отже, щоб визначити
мінімізуючу функцію \texttt{\textbackslash{}overline\{y\}}, треба
розв'язати рівняння Ейлера-Лагранжа.

Отже, ми показали, що якщо функція \texttt{y(x)} мінімізує функціонал
\texttt{J{[}y{]}\ =\ \textbackslash{}int\_a\^{}b\ F(x,y,y\textquotesingle{})\textbackslash{},dx}
у класі гладких функцій із граничними умовами \texttt{y(a)=A} і
\texttt{y(b)=B}, то вона повинна задовольняти рівнянню

\begin{verbatim}
\frac{\partial F}{\partial y} - \frac{d}{dx}\left[\frac{\partial F}{\partial y'}\right] = 0.
\end{verbatim}

(Для спрощення позначень ми прибрали риску над
\texttt{\textbackslash{}overline\{y\}}.) Для ілюстрації теорії
розглянемо приклад.

Пошук мінімальної функціональності
\texttt{J{[}y{]}\ =\ \textbackslash{}int\_\{0\}\^{}\{1\}\ {[}y\^{}2\ +\ y\textquotesingle{}\^{}2{]}\ dx}
.

Спробуємо знайти функцію \texttt{y(x)}, яка проходить через точки (0,0)
і (1,1) і мінімізує функціональність

\begin{verbatim}
P_{0}
\end{verbatim}

З контексту зрозуміло, що шукана функція має бути диференційовною,
оскільки інтегративний вираз залежить від y'. Щоб визначити
\texttt{\textbackslash{}bar\{y\}(x)} , запишемо рівняння
Ейлера--Лагранжа (разом із граничною умовою \texttt{y(0)}=0 та
\texttt{y(1)}=1).

\begin{verbatim}
,  (x,  y, z) \in R^3  ,  u = 0   u_t =
\end{verbatim}

\texttt{y(0)\ =\ 0,} \texttt{y(1)\ =\ 1.} Але
\texttt{F(x,\ y,\ y\textquotesingle{})\ =\ y\^{}2\ +\ y\textquotesingle{}\^{}2}
і, отже,

\begin{enumerate}
\def\labelenumi{\arabic{enumi}.}
\setcounter{enumi}{3}
\tightlist
\item
  Решите аналогичную (см. задачу 3) задачу для двумерного уравнения
\end{enumerate}

(УЧП)

(Продифференцировали \texttt{F} по \texttt{y} и
\texttt{y\textquotesingle{}.} )

Уравнение Эйлера --- Лагранжа принимает вид

Зверніть увагу, що в перших двох задачах коефіцієнти є сталими. Читач
повинен розуміти, що параметр b має бути достатньо малим, а необережний
ряд може відрізнятися

\section{ЗАДАЧИ}\label{ux437ux430ux434ux430ux447ux438}

\begin{verbatim}
- 1. Подставьте разложение (46.8) в задачу (46.7) и получите последовательность задач P_0,\ P_1,\ P_2,\ \dots . 2. Покажите, что нелинейную задачу
\end{verbatim}

откуда

(НУ)

Розв'язуючи це просте диференціальне рівняння з граничними умовами
\texttt{y(0)}=0 та \texttt{y(1)}=1, отримуємо

\begin{verbatim}
u(r, \theta) = u_0(r, \theta) + \frac{1}{4}u_1(r, \theta)
\end{verbatim}

.

Графік цієї функції показано на рис. 44.4. Очевидно, що будь-яка інша
гладка функція, яка задовольняє ті ж граничні умови, надасть більше
значення функції \texttt{J(y)}.

\pandocbounded{\includegraphics[keepaspectratio]{images/_page_332_Figure_5.jpeg}}

Фіг. 44 4. Усі допустимі гладкі криві задовольняють граничні умови
\texttt{y(0)} = 0 та \texttt{y(1)} = 1: a --- допустима крива;
\texttt{\textbackslash{}delta} є мінімізуючою функцією
\texttt{y(x)\ =\ 0.42e\^{}x\ -\ 0.42e\^{}\{-x\}} . \# ПОВАГА 1. Рівняння
Ейлера--Лагранжа аналогічне рівності нульової похідної в
диференціальному численні. Якщо читач пам'ятає, то рівність нульової
похідної не є достатньою умовою екстремуму. Наприклад, похідна функції
\texttt{f(x)\ =\ x\^{}3} в x = 0 дорівнює нулю, але ця точка не є ні np,
ні мінімальною точкою, ні максимальною точкою. Те саме стосується
рівняння Ейлера-Лагранжа: це лише необхідна умова екстремуму, але
недостатня. Це означає, що екстремум можна досягти й на інших функціях.
Однак дуже часто розв'язок рівняння Ейлера-Лагранжа фактично дає
локальний або навіть глобальний мінімум.

\begin{enumerate}
\def\labelenumi{\arabic{enumi}.}
\setcounter{enumi}{1}
\tightlist
\item
  Если у (х) минимизирует функционал
\end{enumerate}

\begin{verbatim}
\Delta u = 0,
\end{verbatim}

тоді площа поверхні, утворена обертанням кривої \texttt{y(x)} навколо
осі x, буде мінімальною (рис. 44.5). Розв'язок рівняння Ейлера-Лагранжа
в цьому випадку буде

цепная линия

\begin{verbatim}
q_{xx} + \varphi_{uu} = 0
\end{verbatim}

де константи \texttt{\textbackslash{}alpha} і
\texttt{\textbackslash{}beta} визначаються з умов \texttt{y(a)} = A та
\texttt{y(b)} = B (тобто ланцюгова лінія має проходити через задані
точки). У цьому випадку досить складно розв'язати рівняння
Ейлера-Лагранжа, тому ми рекомендуємо читачу лише перевірити, чи
задовольняє його лінія ланцюга.

\begin{enumerate}
\def\labelenumi{\arabic{enumi}.}
\setcounter{enumi}{2}
\tightlist
\item
  У цій лекції ми виявили, що
\end{enumerate}

\begin{verbatim}
q_{uu} + q_{vv} = 0
\end{verbatim}

доставляет минимум функционалу Ф і

\pandocbounded{\includegraphics[keepaspectratio]{images/_page_333_Picture_8.jpeg}}

г. 44 5. Мінімальна поверхня обертання.

\begin{enumerate}
\def\labelenumi{\arabic{enumi}.}
\setcounter{enumi}{3}
\tightlist
\item
  Основні закони фізики найчастіше формулюються мовою варіаційних
  принципів, а не диференціальних рівнянь. Прикладами є принцип Ферма
  (світло, поширюючись з однієї точки в іншу, обирає шлях, до якого
  відповідає найкоротшому часу поширення) або принцип Гамільтона (у
  консервативному полі частинка рухається так, що інтеграл дії
\end{enumerate}

\begin{verbatim}
w = f(z)
\end{verbatim}

(кинетическая энергия--- потенциальная энергия) \texttt{dt} буде
мінімальним). Таким чином, фізичні явища розвиваються лише так, щоб ці
функціонали отримували мінімальні значення.

Основні ідеї калькулюсу варіацій можна розширити на функціонали, які
залежать від функції кількох незалежних змінних вигляду П П

\pandocbounded{\includegraphics[keepaspectratio]{images/_page_335_Picture_2.jpeg}}

ринцип Гамільтона стверджує, що матеріальна точка рухатиметься так, що
інтеграл дії буде стаціонарним. Звідси отримують диференціальне рівняння
руху матеріальної точки.

\begin{enumerate}
\def\labelenumi{\arabic{enumi}.}
\setcounter{enumi}{2}
\tightlist
\item
  Показати, що в класі гладких функцій, які задовольняють граничні умови
  \texttt{y(0)}=0 та \texttt{y(\textbackslash{}pi/2)=1} , є мінімальний
  функціонал
\end{enumerate}

Чи знаєте ви, як виглядає графік цієї функції з сегментом {[}0, 1{]}?
Подумай про це. Зверніть увагу, що синусоїдальне перетворення функції
\texttt{f(x)} визначається лише для додатних цілих чисел n.~Іншими
словами, скінченний синус і косинус перетворюють функції у чисельні
послідовності.

\pandocbounded{\includegraphics[keepaspectratio]{images/_page_343_Figure_2.jpeg}}

\begin{verbatim}
Далсе. Приближенное решение   u_n(x,y) = a_1 \varphi_1 + a_2 \varphi_2 + \dots + a_n \varphi_n  .
\end{verbatim}

В это время можно численно найти эначение функционала
J{[}\texttt{u\_n}{]} . Если необходимо найти значение решения
\texttt{u\_n}(x,y) в некоторой точке, можно написать небольшую программу
эля быполнения отой операции. При необходимости можно повторить
вычисления по этой программе с другими значениями п Рис. 452. Блок-схема
метода Ритца.

\section{РОЗВ'ЯЗАННЯ РІВНЯНЬ З ЧАСТКОВИМИ ПОХІДНИМИ МЕТОДАМИ ТЕОРІЇ
ЗБУРЕНЬ}\label{ux440ux43eux437ux432ux44fux437ux430ux43dux43dux44f-ux440ux456ux432ux43dux44fux43dux44c-ux437-ux447ux430ux441ux442ux43aux43eux432ux438ux43cux438-ux43fux43eux445ux456ux434ux43dux438ux43cux438-ux43cux435ux442ux43eux434ux430ux43cux438-ux442ux435ux43eux440ux456ux457-ux437ux431ux443ux440ux435ux43dux44c}

Мета лекції - показати, як складні задачі з нелінійностями, змінними
коефіцієнтами або нерегулярними областями можна розв'язувати як збурення
простіших задач.

Наприклад, можна розглядати задачу Діріхле в деформованій області

\begin{verbatim}
\Delta u = 0, \qquad 0 < r < 1 + \frac{1}{4}\sin\theta,
\end{verbatim}

з граничною умовою

\begin{verbatim}
u\!\left(1+\frac{1}{4}\sin\theta,\theta\right)=\cos\theta, \qquad 0 \le \theta \le 2\pi.
\end{verbatim}

Очень часто задачу, решение которой не известно, можно путем
непрерывного изменения каких-то параметров свести к близкой, но легко
решаемой задаче. Тогда и решение исходной задачи будет близко решению
модифицирующей задачи.

Например, уравнение Лапласа

\$\$
\texttt{\textbackslash{}begin\{cases\}\ 1,\ \&\ n\ =\ 1,\ \textbackslash{}\textbackslash{}\ 0,\ \&\ n\ =\ 2,\ 3,\ \textbackslash{}dots,\ \textbackslash{}end\{cases\}}
за допомогою сімейства задач, параметризованих числом
\texttt{\textbackslash{}varepsilon}, де
\texttt{0\ \textbackslash{}le\ \textbackslash{}varepsilon\ \textbackslash{}le\ 1}.

\texttt{Розглядається\ область\ \textquotesingle{}0\ \textless{}\ r\ \textless{}\ 1\textquotesingle{},\ \textquotesingle{}0\ \textbackslash{}le\ \textbackslash{}theta\ \textless{}\ 2\textbackslash{}pi\textquotesingle{}\ з\ граничною\ умовою\ \textquotesingle{}u(1,\textbackslash{}theta)=g(\textbackslash{}theta)\textquotesingle{}.}

за деяких граничних умов (наприклад,
\texttt{g(\textbackslash{}theta)\ =\ \textbackslash{}cos\ \textbackslash{}theta}
) має два різних матеріальні розв'язки, тоді як за інших (наприклад,
\texttt{g(\textbackslash{}theta)\ =\ A\ \textbackslash{}cos\ \textbackslash{}theta}
з A \textgreater{} 20.65) задача не має реального розв'язку. --- Ед. Ед.

4) Краевая задача

\texttt{можно\ непрерывно\ модифицировать\ в\ нелинейное\ уравнение\ \textquotesingle{}\textbackslash{}Delta\ u\ =\ u\^{}2\ =\ 0\textquotesingle{}\ 1).}

Фіг. 46.1 показує схему розв'язання нелінійної задачі шляхом її зведення
до розв'язку рівнянняЛапласа 2)

\end{document}
